\documentclass[a4paper, 	% Seitenformat
		12pt, 								% Schriftgröße
		bibliography=totoc, 	% Literaturverzeichnis in das Inhaltsverzeichnis
		index=totoc, 	% Index in das Inhaltsverzeichnis
		abstracton, 	% mit Abstrakt
		headsepline, 	% Trennlinie für die Kopfzeile
		%footnosepline,	% Trennlinie für die Fusszeile
		]{scrreprt}

\usepackage{latexsym}		      % Latex-Symbole
\usepackage[latin1]{inputenc}	% Eingabekodierungen
\usepackage[ngerman]{babel}	  % Mehrsprachenumgebung
%\usepackage{ulem}
% Bilder
%\usepackage{pst-all}		% Zeichnungen in Latex (kein pdflatex)
\usepackage{subfig}	  	% Unterteilungen von Abbildungen
\usepackage{float}		  % Gleitobjekte
\usepackage{graphicx}		% Einbinden von Graphiken aus Dateien
\usepackage{rotating}   % Die folgenden Packages sind meist hilfreich für das Erstellen von Arbeiten
\usepackage{mathrsfs}
\usepackage{eurosym,bm,amsmath,amssymb,verbatim}
\usepackage{listings}
\usepackage[T1]{fontenc}
\usepackage{setspace}


%------------------------------------------------------------------------

\parindent0em 			% kein Einzug nach einer Leerzeile
\topmargin-20.4mm   % Für die Ränder
\textheight245mm

\textwidth157mm \oddsidemargin11mm \evensidemargin11mm

%%%%%%%%%%%%%%%%%%%%%%%%%%%%%%%%%%%%%%%%%%%%%%%%%%%%%%%%%%%%%%%%%%%%%%%
%%%%%%%%%%%%%%%%%%%%%%%%%%%%%%%%%%%%%%%%%%%%%%%%%%%%%%%%%%%%%%%%%%%%%%%
%%%%%%%%%%%%%%%%%%%%%%%%%%%%%%%%%%%%%%%%%%%%%%%%%%%%%%%%%%%%%%%%%%%%%%%

\begin{document}					% Start des Dokuments

\begin{titlepage}

\hspace{12cm}\scalebox{0.2}{\includegraphics{images/uni_aug_logo.png}}

\vspace*{1cm}
\large\textbf{Thema der wissenschaftlichen Arbeit}\\[0.5cm]

\normalsize Freie wissenschaftliche Arbeit\\

zur Erlangung des akademischen Grades\\

"`Bachelor of Arts"' etc.\\

Studiengang: Betriebswirtschaftslehre etc.\\

\textbf{an der Wirtschaftswissenschaftlichen Fakult\"at der Universit\"at \mbox{Augsburg}}\\[0.8cm]

Lehrstuhl f\"ur Statistik und Data Science\\[1.5cm]
%
%
\noindent\begin{tabular}[h]{@{}ll}
%\begin{tabular}{l l}
Eingereicht bei: & \\
\\
Betreuer: 			 & Betreuer   \\
\\
Vorgelegt von:   &Verfasser               \\
\\
Adresse:         &Stra\ss e                  \\[0.2cm]
                 &Ort                     \\
\\
Matrikel-Nr.:    &1234567                 \\
\\
Email:           &Verfasser@email.de      \\
\\
\end{tabular}

\vfill
Augsburg, im September 2024

\end{titlepage}


\newpage
\pagenumbering{Roman}


%\tableofcontents		    % Inhaltsverzeichnis
%\newpage
%\listoffigures			    % Abbildungsverzeichnis
%\listoftables				  % Tabellenverzeichnis

\clearpage
\pagenumbering{arabic}
\onehalfspacing 			     % Für richtigen Zeilenabstand

%%%%%%%%%%%%%%%%%%%%%%%%%%%%%%%%%%%%%%%%%%%%%%%%%%%%%%%%%%%%%%%%%%%%%%%
%%%%%%%%%%%%%%%%%%%%%%%%%%%%%%%%%%%%%%%%%%%%%%%%%%%%%%%%%%%%%%%%%%%%%%%
%%%%%%%%%%%%%%%%%%%%%%%%%%%%%%%%%%%%%%%%%%%%%%%%%%%%%%%%%%%%%%%%%%%%%%%

\input{Kapitel/kapitel1}  % Zur besseren Strukturierung empfiehlt es sich das Dokument aufzuteilen.
%\newpage
%\input{Kapitel/kapitel2}
%\newpage
%\input{Kapitel/kapitel3}
%\newpage
%\input{Kapitel/kapitel4}
%\newpage
%\input{Kapitel/kapitel5}
%\vspace{5em}


\clearpage

\bibliographystyle{acm}
\bibliography{Literatur}	           % Literaturverzeichnis liegt in der Datei Literatur

\end{document}											 % Ende des Dokuments
